\section{System Management}
\label{sec:Management}

This chapter covers system management tasks: creating and deleting models,
compactions, exports, imports, rebuilds, tags and integrity checks.

\subsection{Model lifecycle}

\subsubsection{Creating a model}

Models must be explicitly registered before clients can use them:

\begin{lstlisting}[language=bash,caption={Register a model}]
curl -s -X POST http://localhost:8081/admin/models/demo | jq .
\end{lstlisting}

Optionally provide a display name:

\begin{lstlisting}[language=bash,caption={Register with name}]
curl -s -X POST http://localhost:8081/admin/models/demo \
  -H "Content-Type: application/json" \
  -d '{"modelName":"Demo Model"}' | jq .
\end{lstlisting}

The creating user (resolved from the current identity) is seeded as the
initial model admin, writer and reader in the ACL.

\subsubsection{Deleting a model}

\begin{lstlisting}[language=bash,caption={Delete a model}]
curl -s -X DELETE http://localhost:8081/admin/models/demo | jq .
\end{lstlisting}

Deletion is refused while the model has active sessions. To override:

\begin{lstlisting}[language=bash,caption={Force delete}]
curl -s -X DELETE "http://localhost:8081/admin/models/demo?force=true" | jq .
\end{lstlisting}

\subsection{Compaction}

Compaction prunes old op-log entries and merge metadata below a configurable
retention watermark to bound storage growth.

\subsubsection{Running compaction}

\begin{lstlisting}[language=bash,caption={Compact a model}]
curl -s -X POST http://localhost:8081/admin/models/demo/compact | jq .
\end{lstlisting}

The response reports:

\begin{itemize}
\item \texttt{modelId}
\item \texttt{headRevision} --- current model head
\item \texttt{committedHorizon} --- revision below which compaction can safely operate
\item \texttt{retainRevisions} --- how many recent revisions are preserved
\item \texttt{watermarkRevision} --- revisions below this are eligible for pruning
\item \texttt{commitsRemoved} --- number of commit nodes deleted
\item \texttt{opsRemoved} --- number of op nodes deleted
\item \texttt{propertyClocksRemoved} --- number of orphan PropertyClock
  nodes removed
\item \texttt{tombstonesEligible} --- number of tombstones eligible under
  watermark but retained for delete safety
\item \texttt{clockFieldsEligible} --- number of field clock metadata
  entries eligible under watermark but retained
\end{itemize}

\subsubsection{Compaction safety}

Compaction is designed to be safe by default:

\begin{itemize}
\item Only commits and ops with revision below the watermark are removed.
\item The watermark is derived from the committed revision horizon
  (\texttt{latestCommitRevision - retainRevisions}).
\item Tombstones are \emph{never} removed by compaction. They are reported
  as eligible but retained to prevent stale resurrection.
\item Field clocks (\texttt{*\_lamport}, \texttt{*\_clientId}) on
  materialized nodes are retained.
\item Only orphan \texttt{PropertyClock} nodes (no longer referenced by any
  materialized entity) below the watermark are pruned.
\item Compaction never allows a stale \texttt{Create*} or \texttt{Update*}
  to beat a logically newer delete.
\end{itemize}

\subsection{Export}

Export produces a self-contained package for one model timeline, suitable for
backup or migration.

\begin{lstlisting}[language=bash,caption={Export a model}]
curl -s http://localhost:8081/admin/models/demo/export | jq .
\end{lstlisting}

The export package contains:

\begin{description}
\item[\texttt{format}] Package format identifier (\texttt{archi-model-export-v1}).
\item[\texttt{exportedAt}] UTC timestamp.
\item[\texttt{model}] Model catalog metadata (modelId, modelName, headRevision).
\item[\texttt{accessControl}] Per-model ACL snapshot (optional).
\item[\texttt{snapshot}] Materialized current state at HEAD. Contains
  folders, elements, relationships, views, view objects, connections,
  folder memberships and view object child members.
\item[\texttt{opBatches}] Ordered committed op-log history. Each batch
  contains modelId, opBatchId, baseRevision, assignedRevisionRange,
  timestamp and the ops array with normalized causal metadata.
\item[\texttt{tags}] Immutable tags with their historical snapshots.
\end{description}

The package format is documented in detail in \filepath{server/EXPORT\_FORMAT.md}.

\subsubsection{What export preserves}

\begin{itemize}
\item Model id and name
\item Head revision
\item Full committed op-log history
\item Current materialized state
\item Folders and folder memberships
\item Views, view objects, connections and notation payloads
\item Immutable tags and their historical snapshots
\item Model-scoped ACLs
\end{itemize}

\subsection{Import}

Import restores a model from an export package.

\begin{lstlisting}[language=bash,caption={Import a model}]
curl -s -X POST http://localhost:8081/admin/models/import \
  -H "Content-Type: application/json" \
  -d @demo-export.json | jq .
\end{lstlisting}

\subsubsection{Import behavior}

\begin{itemize}
\item The \texttt{format} field must be \texttt{archi-model-export-v1}.
\item \texttt{model.modelId} must be present.
\item Existing models are rejected unless \texttt{?overwrite=true}.
\item Overwrite is rejected if the target model has active sessions.
\item If overwriting, the existing model is deleted first.
\item ACLs are restored if present in the package.
\item Op batches are restored first, then head revision is set from the
  imported history.
\item Tags are restored after the model timeline.
\end{itemize}

\subsubsection{Overwrite with import}

\begin{lstlisting}[language=bash,caption={Import with overwrite}]
curl -s -X POST "http://localhost:8081/admin/models/import?overwrite=true" \
  -H "Content-Type: application/json" \
  -d @demo-export.json | jq .
\end{lstlisting}

\subsubsection{Import requirements}

For non-empty models, the package must include \texttt{opBatches}. Import
rejects packages that have a non-zero head revision, non-empty materialized
content or historical tags but no op-log history. This prevents restoring
current state while silently losing the committed revision history.

\subsection{Rebuild}

Rebuild clears the materialized state for a model and replays the entire
op-log to reconstruct it. This is useful after corruption or for verifying
that the op-log faithfully reproduces the materialized state.

\begin{lstlisting}[language=bash,caption={Rebuild a model}]
curl -s -X POST http://localhost:8081/models/demo/rebuild | jq .
\end{lstlisting}

Or via the admin endpoint:

\begin{lstlisting}[language=bash,caption={Rebuild and get status}]
curl -s -X POST http://localhost:8081/admin/models/demo/rebuild-and-status | jq .
\end{lstlisting}

Rebuild behavior:

\begin{enumerate}
\item Clears all materialized nodes for the model (elements, relationships,
  views, view objects, connections, folders and metadata).
\item Recreates built-in root folders.
\item Replays all committed ops from the op-log in revision order.
\item Updates head revision to match the highest committed revision.
\item Tags are unaffected (they are stored separately).
\end{enumerate}

\subsection{Tags}

Tags are immutable named pointers to historical revisions. They let you
capture known-good states for later reference or checkout.

\subsubsection{Creating a tag}

\begin{lstlisting}[language=bash,caption={Create a tag}]
curl -s -X POST http://localhost:8081/admin/models/demo/tags \
  -H "Content-Type: application/json" \
  -d '{"tagName":"v1.2","description":"Sprint 12 baseline"}' | jq .
\end{lstlisting}

The tag captures the current HEAD revision and stores a snapshot of the
materialized state at that point.

\subsubsection{Reading at a tag}

Clients can join at a tagged revision by specifying the tag name as the
\texttt{ref} in the \texttt{Join} message. The snapshot endpoint also accepts
\texttt{?ref=<tagName>}:

\begin{lstlisting}[language=bash,caption={Read snapshot at tag}]
curl -s "http://localhost:8081/models/demo/snapshot?ref=v1.2" | jq .
\end{lstlisting}

Tagged refs are read-only. Writes, locks and presence are only allowed on
\texttt{HEAD}.

\subsubsection{Tag naming}

Use descriptive names that remain meaningful outside the codebase:

\begin{itemize}
\item \texttt{v1.2} --- version numbers
\item \texttt{milestone-3} --- project milestones
\item \texttt{approved-2026-03-08} --- date-stamped approvals
\end{itemize}

Tag deletion is disabled by default. Re-enable it only with
\texttt{app.tags.allow-delete=true} if you explicitly want mutable tag
administration.

\subsection{Integrity checks}

The integrity endpoint checks for structural problems in the materialized
graph:

\begin{lstlisting}[language=bash,caption={Check integrity}]
curl -s http://localhost:8081/admin/models/demo/integrity | jq .
\end{lstlisting}

The response reports:

\begin{itemize}
\item Missing references: entities that reference ids not present in the
  graph (e.g. a view object whose \texttt{representsId} points to a deleted
  element).
\item Orphan nodes: materialized nodes not reachable from the model root.
\item Dangling relationships: relationships where source or target element
  is missing.
\end{itemize}

Integrity issues indicate either a bug in the apply path or manual Neo4j
modifications. If issues are found, run a rebuild to reconstruct the
materialized state from the op-log.

\subsection{Consistency checks}

Optional post-commit consistency checks can be enabled with
\texttt{ARCHIMESH\_CONSISTENCY\_CHECK=true}. When enabled, the server runs
a lightweight consistency verification after each \texttt{SubmitOps}
commit, comparing the materialized state against the expected state derived
from the op-log. Failures are logged at WARN level.

This is disabled by default because it adds latency to every submit.
Enable it during debugging or after infrastructure changes.

\subsection{ACL management}

Per-model ACLs restrict model access to specific users:

\begin{lstlisting}[language=bash,caption={Read ACL}]
curl -s http://localhost:8081/admin/models/demo/acl | jq .
\end{lstlisting}

\begin{lstlisting}[language=bash,caption={Update ACL}]
curl -s -X PUT http://localhost:8081/admin/models/demo/acl \
  -H "Content-Type: application/json" \
  -d '{"adminUsers":["alice"],"writerUsers":["alice","bob"],"readerUsers":["carol"]}' | jq .
\end{lstlisting}

When a model has a configured ACL, it overrides generic reader/writer role
checks for that model. Catalog-wide admin actions still require the global
admin role.

\subsection{Status monitoring}

Check model health at a glance:

\begin{lstlisting}[language=bash,caption={Model status}]
curl -s http://localhost:8081/admin/models/demo/status | jq .
\end{lstlisting}

The response includes:

\begin{itemize}
\item \texttt{modelId}, \texttt{modelName}
\item \texttt{headRevision}
\item Element, relationship, view, folder, view object and connection counts
\item Consistency check results
\item Compaction status (last compaction timestamp, watermark, commits removed)
\item Active session count
\end{itemize}
